\begin{refsection}

\begin{center}
{\Large Приложение Б. Обоснование методики вычислительных экспериментов и статистической обработки результатов для задачи mTSP}

\vspace{0.5em}
{\large к проекту \texttt{mtsp-routing-optimization}}
\end{center}

\vspace{1em}

\section{Назначение приложения и его границы}

Настоящее приложение фиксирует методику постановки вычислительных экспериментов для задачи множественного коммивояжёра и правила статистической обработки получаемых результатов. Его цель состоит не в повторном описании реализованных эвристик, а в формальном обосновании того, как именно следует строить тестовый набор, сколько независимых наблюдений требуется для сравнения алгоритмов и каким образом интерпретировать различия между ними при подготовке итогового текста курсовой работы.

Фокус данного приложения ограничен экспериментальным дизайном. Вопросы корректности построения маршрутов, допустимости локальных преобразований и структурных причин, по которым более сильные эвристики превосходят базовые схемы, вынесены в соседнее математическое приложение. Такое разделение соответствует общим рекомендациям по экспериментальному анализу алгоритмов: описание алгоритма, экспериментальный протокол и статистическая интерпретация результатов должны быть разведены по уровням, чтобы не смешивать свойства метода с качеством постановки вычислительного эксперимента \cite{johnson1999,hooker1995}.

При подготовке настоящего приложения использованы локальные артефакты проекта:
\begin{enumerate}[label=\arabic*.]
  \item конфигурация основного benchmark-контура \texttt{experiments/config.json};
  \item текущие результаты основного прогона \texttt{data/results/mtsp\_results.csv};
  \item сводка сравнения с внешним reference baseline \texttt{data/results/mtsp\_reference\_summary.csv};
  \item парные статистики по текущему uniform-набору \texttt{data/results/methodology\_pairwise\_current.csv};
  \item pilot-оценки разброса по различным значениям \texttt{seed} из \texttt{data/results/methodology\_pilot\_summary.csv} и \texttt{data/results/methodology\_pilot\_ci\_projection.csv}.
\end{enumerate}

Поэтому далее последовательно различаются два уровня утверждений:
\begin{enumerate}[label=\arabic*.]
  \item \textbf{подтверждённые факты} --- выводы, которые уже следуют из текущих CSV-артефактов и pilot-расчётов;
  \item \textbf{зафиксированный итоговый протокол} --- рекомендуемая схема финальных экспериментов для курсовой, которая частично уже отражена в конфигурации и генераторе, но ещё не полностью прогнана в основном наборе результатов.
\end{enumerate}

\section{Что уже подтверждено текущими артефактами}

Текущий основной CSV \texttt{mtsp\_results.csv} содержит \(480\) строк, то есть результаты \(4\) solver'ов на \(120\) различных экземплярах. Эти данные покрывают размерности \(n \in \{20,50,100,200\}\) и числа коммивояжёров \(m \in \{2,3,5\}\), но относятся только к uniform-инстансам: в таблице отсутствует признак семейства инстансов, а пути к файлам указывают на \texttt{data/mtsp/generated/n*\_m*\_r*.txt}. Следовательно, текущий основной CSV пока подтверждает только предварительное сравнение алгоритмов на однородной uniform-сетке.

С точки зрения качества решений текущая основная таблица уже содержит полезный сигнал. По сводке \texttt{mtsp\_reference\_summary.csv} средний относительный разрыв к reference baseline на базе OR-Tools составляет \(14.477\%\) для \texttt{lkh-wrapper-v2}, \(20.769\%\) для \texttt{grasp}, \(47.413\%\) для \texttt{2opt+greed} и \(114.730\%\) для \texttt{rand+nn}. Следовательно, даже ограниченный uniform-набор уже подтверждает, что \texttt{lkh-wrapper-v2} является сильнейшим из текущих собственных solver'ов по средней близости к внешнему эталону.

Дополнительная проверка по \texttt{methodology\_pairwise\_current.csv} показывает, что при сравнении \texttt{grasp} и \texttt{lkh-wrapper-v2} на одних и тех же uniform-экземплярах \texttt{lkh-wrapper-v2} имеет лучшее среднее значение целевой функции в \(10\) из \(12\) страт \((n,m)\). При этом в \(8\) из \(12\) страт \(95\%\)-й доверительный интервал для средней парной разности не содержит нуль, а в \(4\) стратах остаётся неоднозначность, требующая подтверждающего объёма данных. Наибольшая текущая полуширина такого интервала при \(N=10\) достигает \(39.612\), что выше порога \(h_D \le 35\), выбранного далее в качестве базового критерия точности. Это означает, что \(N=10\) в текущем проекте уже достаточно для screening-сравнения, но ещё не гарантирует одинаково узкий интервал для всех страт.

Отдельный pilot по различным значениям \texttt{seed} уже построен на пяти экземплярах: \texttt{uniform\_n100\_m3\_pilot.txt}, \texttt{uniform\_n200\_m5\_pilot.txt}, \texttt{clustered\_center\_n100\_m3\_pilot.txt}, \texttt{clustered\_offset\_depot\_n100\_m3\_pilot.txt} и \texttt{mixed\_outliers\_n100\_m3\_pilot.txt}. Для каждого из трёх стохастических solver'ов выполнено по \(20\) независимых запусков, что позволяет оценить внутрисубъектную вариативность и проектировать разумный выбор \(R\) для основной сетки. Наиболее важен worst-case для \texttt{grasp}: по \texttt{methodology\_pilot\_ci\_projection.csv} максимальная относительная полуширина \(95\%\)-го доверительного интервала по seed составляет \(0.917\%\) при \(R=20\), \(0.731\%\) при \(R=30\) и \(0.557\%\) при \(R=50\). Эти значения непосредственно используются ниже как опора для выбора базового, подтверждающего и строгого режимов.

\begin{table}[h]
\centering
\caption{Что уже подтверждено текущими артефактами}
\begin{tabular}{p{0.29\linewidth}p{0.33\linewidth}p{0.27\linewidth}}
\toprule
Артефакт & Подтверждённый вывод & Ограничение \\
\midrule
\texttt{mtsp\_results.csv} & \(120\) экземпляров, \(4\) solver'а, uniform-сетка по \(n \in \{20,50,100,200\}\), \(m \in \{2,3,5\}\) & Нет семейства инстансов, нет режима \(m=7\), стохастические solver'ы запущены с фиксированным \texttt{seed=42} \\
\texttt{mtsp\_reference\_summary.csv} & Средний gap к OR-Tools минимален у \texttt{lkh-wrapper-v2} & Сравнение выполняется с фиксированным внешним бюджетом \(5\) с на экземпляр, а не в matched-budget постановке \\
\texttt{methodology\_pairwise\_current.csv} & \texttt{lkh-wrapper-v2} лучше \texttt{grasp} в \(10/12\) uniform-страт; в \(8/12\) CI исключает нуль & При \(N=10\) часть страт остаётся статистически спорной \\
\texttt{methodology\_pilot\_*} & Pilot по \(20\) seed даёт оценку разброса для трёх стохастических solver'ов на разных семействах инстансов & Pilot покрывает только \(5\) специально выбранных экземпляров, а не весь основной грид \\
\bottomrule
\end{tabular}
\end{table}

\section{Формальная статистическая модель эксперимента}

В проекте рассматривается задача mTSP с фиксированным числом коммивояжёров и целевой функцией \texttt{MINSUM}; обозначения постановки задачи вынесены в отдельное математическое приложение. Используемая здесь терминология согласуется со стандартной постановкой задачи mTSP \cite{bektas2006}. Для целей экспериментального дизайна достаточно считать, что каждая страта задаётся тройкой
\[
s = (family, n, m),
\]
где \(family\) обозначает семейство инстансов, \(n\) --- число вершин, а \(m\) --- число маршрутов.

Пусть \(a\) обозначает алгоритм, \(j \in \{1,\dots,N_s\}\) --- номер экземпляра внутри фиксированной страты \(s\), а \(r \in \{1,\dots,R_a\}\) --- номер независимого запуска на одном и том же экземпляре. Для детерминированных solver'ов принимается \(R_a = 1\); для стохастических solver'ов \(R_a = R\), где \(R\) выбирается по протоколу, зафиксированному в разделе~5. Результат запуска по основной метрике качества обозначим через
\[
X_{a,j,r}.
\]

Для фиксированных \(a\) и \(j\) выборочное среднее и выборочное стандартное отклонение по повторным seed-запускам равны
\[
\bar X_{a,j} = \frac{1}{R}\sum_{r=1}^{R} X_{a,j,r},
\qquad
s_{a,j} = \sqrt{\frac{1}{R-1}\sum_{r=1}^{R}\left(X_{a,j,r}-\bar X_{a,j}\right)^2}.
\]
Именно величина \(\bar X_{a,j}\) используется далее как оценка среднего поведения стохастического solver'а на фиксированном экземпляре, а \(s_{a,j}\) --- как оценка внутрисубъектной вариативности. При достаточном числе повторов доверительный интервал для среднего строится по стандартной \(t\)-схеме \cite{nistNormal,nistMeanCI}:
\[
\bar X_{a,j} \pm t_{0.975,R-1}\frac{s_{a,j}}{\sqrt{R}}.
\]

Для сравнения двух алгоритмов \(a\) и \(b\) внутри одной страты необходимо использовать парную постановку, поскольку оба алгоритма запускаются на одном и том же наборе экземпляров. Для каждого экземпляра \(j\) вводится разность
\[
D_j = \bar X_{a,j} - \bar X_{b,j}.
\]
Тогда вся задача сравнения сводится к анализу выборки \(D_1,\dots,D_{N_s}\), а доверительный интервал для средней разности имеет вид \cite{psuPaired,nistMeanCI}
\[
\bar D \pm t_{0.975,N_s-1}\frac{s_D}{\sqrt{N_s}},
\]
где
\[
\bar D = \frac{1}{N_s}\sum_{j=1}^{N_s} D_j,
\qquad
s_D = \sqrt{\frac{1}{N_s-1}\sum_{j=1}^{N_s}(D_j - \bar D)^2}.
\]

Такая blocked/paired-схема принципиальна: она вычитает вариативность, связанную с самими экземплярами, и тем самым делает сравнение мощнее, чем независимое сопоставление средних \cite{psuPaired,johnson1999}. Основной количественный вывод в настоящем приложении опирается именно на среднюю парную разность и её \(95\%\)-й доверительный интервал.

Для проверки гипотезы о нулевой средней разности используется парный \(t\)-критерий. В качестве непараметрической проверки устойчивости вывода дополнительно применяется критерий Уилкоксона для signed-rank разностей \cite{scipyWilcoxon}. Если в одной страте выполняется несколько парных сравнений между solver'ами, семейство \(p\)-значений корректируется процедурой Холма \cite{holm1979}. Это правило не заменяет доверительные интервалы, а лишь ограничивает риск ложных срабатываний при множественной проверке.

Для стохастических solver'ов удобно контролировать не только абсолютную, но и относительную точность оценки среднего по повторам. Поэтому вводится относительная полуширина доверительного интервала
\[
q_{a,j} =
100\% \cdot
\frac{t_{0.975,R-1}\, s_{a,j}}{\sqrt{R}\,\bar X_{a,j}}.
\]
В дальнейшем пороги на \(q\) используются как операционное правило выбора \(R\). Аналогично, для парного сравнения в страте контролируется абсолютная полуширина
\[
h_D = t_{0.975,N_s-1}\frac{s_D}{\sqrt{N_s}},
\]
которая задаёт точность оценки средней разности между двумя алгоритмами.

Если задача состоит не только в построении доверительного интервала, но и в гарантированном обнаружении practically important difference \(\Delta^\*\), то для paired \(t\)-design используется стандартизованный эффект
\[
d = \frac{\Delta^\*}{\sigma_D},
\]
а расчёт требуемого \(N_s\) опирается на стандартные power-функции для one-sample / paired \(t\)-test \cite{statsmodelsTTestPower}. В настоящем приложении этот подход используется не для автоматической подгонки каждого эксперимента, а для фиксации реалистичных порогов основного и подтверждающего режимов.

\section{Аудит текущего состояния экспериментального контура}

На момент подготовки настоящего приложения экспериментальный контур проекта уже частично приведён к итоговой схеме, но ещё не завершён. На уровне конфигурации \texttt{experiments/config.json} уже зафиксированы четыре семейства инстансов --- \texttt{uniform}, \texttt{clustered-center}, \texttt{clustered-offset-depot}, \texttt{mixed-outliers} --- и дополнительный stress-режим \(m=7\) для \(n \in \{100,200\}\). Кроме того, основной solver в конфигурации уже переведён на \texttt{lkh-wrapper-v2}. Однако фактические артефакты основного прогона пока отстают от этой схемы.

Во-первых, каталог \texttt{data/mtsp/generated} на текущем этапе содержит только \(120\) uniform-файлов без семейств \texttt{clustered-*}, без \texttt{mixed-outliers} и без экземпляров с \(m=7\). Соответственно, и основной CSV отражает только uniform-сетку. Во-вторых, для стохастических solver'ов в основном прогоне по-прежнему используется фиксированный \texttt{seed=42}, то есть в \texttt{mtsp\_results.csv} есть только по одному качественному наблюдению на solver и экземпляр, а не выборка независимых запусков. В-третьих, сравнение с OR-Tools уже присутствует и полезно как внешний ориентир по качеству, но не образует matched-budget протокола, поскольку OR-Tools получает фиксированные \(5\) секунд на экземпляр, тогда как собственные solver'ы в текущем контуре работают существенно быстрее.

\begin{table}[h]
\centering
\caption{Сопоставление целевого дизайна и текущего состояния основного контура}
\begin{tabular}{p{0.31\linewidth}p{0.25\linewidth}p{0.27\linewidth}}
\toprule
Характеристика & Целевой дизайн & Текущее состояние \\
\midrule
Семейства инстансов & \(4\): \texttt{uniform}, \texttt{clustered-center}, \texttt{clustered-offset-depot}, \texttt{mixed-outliers} & В основном гриде отражён только \texttt{uniform}; остальные семейства есть лишь в pilot \\
Размерности \(n\) & \(20, 50, 100, 200\) & Подтверждены в основном CSV \\
Число коммивояжёров \(m\) & \(2,3,5\), а для \(n \in \{100,200\}\) дополнительно \(7\) & В основном CSV присутствуют только \(2,3,5\) \\
Число страт \((family,n,m)\) & \(56\) & В основном CSV фактически \(12\) uniform-страт \\
Число экземпляров на базовом уровне & \(N=10\) на страту, всего \(560\) файлов & В каталоге основного грида сейчас \(120\) uniform-файлов \\
Повторы для стохастических solver'ов & \(R \ge 20\) & В основном CSV один фиксированный \texttt{seed=42}; повторность пока есть только в pilot \\
Сравнение с OR-Tools & Внешний baseline; matched-budget при необходимости ставится отдельно & Есть reference baseline с \(5\) с на экземпляр, matched-budget отдельно не поставлен \\
\bottomrule
\end{tabular}
\end{table}

Из этой таблицы следует два принципиальных вывода. С одной стороны, текущий контур уже достаточно зрелый, чтобы обосновывать выбор solver'ов, строить pilot-оценки разброса и фиксировать итоговый протокол. С другой стороны, итоговые утверждения уровня ``алгоритм \(A\) лучше алгоритма \(B\) на заявленном benchmark-наборе'' пока нельзя переносить на все четыре семейства инстансов, пока основной грид не будет действительно прогнан в полном конфигурационном объёме.

\section{Зафиксированный итоговый протокол для курсовой}

\subsection{Состав тестовых данных}

Итоговый экспериментальный набор для курсовой фиксируется по стратам \((family,n,m)\) и включает:
\begin{enumerate}[label=\arabic*.]
  \item семейства \texttt{uniform}, \texttt{clustered-center}, \texttt{clustered-offset-depot}, \texttt{mixed-outliers};
  \item размерности \(n \in \{20,50,100,200\}\);
  \item числа маршрутов \(m \in \{2,3,5\}\), а для \(n \in \{100,200\}\) дополнительно \(m=7\).
\end{enumerate}

Такой набор образует \(56\) страт: по \(3\) значения \(m\) для \(n=20\) и \(n=50\), и по \(4\) значения \(m\) для \(n=100\) и \(n=200\), умноженные на четыре семейства инстансов. При базовом объёме \(N=10\) это даёт \(560\) различных экземпляров. Именно эта сетка рассматривается как целевой финальный benchmark для курсовой; до момента её полного прогона все обобщения на четыре семейства должны сопровождаться явной оговоркой, что часть утверждений пока подтверждена только uniform-артефактами и отдельным pilot.

\subsection{Число экземпляров, число повторов и формальные пороги}

Базовый, подтверждающий и строгий режимы фиксируются не как свободные рекомендации, а как часть итогового протокола. Их значения сведены в таблицу~\ref{tab:final-protocol}.

\begin{table}[h]
\centering
\caption{Зафиксированные параметры итогового протокола}
\label{tab:final-protocol}
\begin{tabular}{p{0.35\linewidth}ccc}
\toprule
Параметр & Базовый режим & Подтверждающий режим & Строгий локальный режим \\
\midrule
Число экземпляров на страту \(N\) & \(10\) & \(20\) & \(>20\) при необходимости \\
Число повторов по \texttt{seed} для стохастических solver'ов \(R\) & \(20\) & \(30\) & \(50\) \\
Целевая мощность \(1-\beta\) & \(0.80\) & \(0.90\) & не ниже \(0.90\) \\
Practically important difference \(\Delta^\*\) & \(25/35/45\) в абсолютной шкале или \(1\%/1.5\%/2\%\) gap в относительной шкале & те же пороги, но с более узким интервалом & используется только для спорных локальных сравнений \\
Допустимая полуширина по повторам \(q\) & \(q \le 1\%\) & \(q \le 0.75\%\) & стремление к \(q \approx 0.5\%\) \\
Допустимая полуширина по парной разности \(h_D\) & \(h_D \le 35\) & \(h_D \le 25\) & уже после локального усиления \(N\) \\
Множественные сравнения & Holm внутри страты & Holm внутри страты & Holm внутри страты \\
\bottomrule
\end{tabular}
\end{table}

Фиксация именно этих значений согласуется с текущими локальными артефактами. Pilot по \texttt{grasp} показывает, что при \(R=20\) worst-case относительная полуширина уже удерживается на уровне около \(0.917\%\), а переход к \(R=30\) и \(R=50\) снижает её примерно до \(0.731\%\) и \(0.557\%\) соответственно. Поэтому \(R=20\) корректно трактовать как основной рабочий режим для полной сетки, \(R=30\) --- как подтверждающий режим для спорных страт и финальной пары лидеров, а \(R=50\) --- как дорогой локальный режим, а не как обязательный стандарт для всего benchmark-а.

Аналогично, текущие парные uniform-статистики показывают, что при \(N=10\) максимальная полуширина доверительного интервала для средней разности достигает \(39.612\). Следовательно, \(N=10\) следует честно называть screening-уровнем: он уже достаточен для отсева слабых solver'ов и для предварительного ранжирования, но не даёт одинаково узкой оценки во всех стратах. Подтверждающий режим \(N=20\) нужен именно для тех случаев, где базовый объём оставляет значимую неопределённость.

\subsection{Правило эскалации объёма данных}

После завершения базового прогона каждая страта проходит одно и то же правило принятия решения:
\begin{enumerate}[label=\arabic*.]
  \item если \(95\%\)-й доверительный интервал для средней парной разности не содержит нуль и практически значимая граница \(\Delta^\*\) также лежит вне интервала, базовый результат принимается без расширения;
  \item если интервал пересекает нуль или границу practically important difference, страта переводится в подтверждающий режим \(N=20, R=30\);
  \item если у лидирующего стохастического solver'а нарушается порог \(q \le 1\%\) в базовом режиме или \(q \le 0.75\%\) в подтверждающем режиме, число повторов увеличивается;
  \item если для итогового сравнения двух лучших solver'ов требуется особенно узкая оценка среднего поведения, используется локальный режим \(R=50\) и, при необходимости, \(N>20\).
\end{enumerate}

Такое правило делает протокол decision-complete: оно заранее определяет, когда базового объёма достаточно, а когда требуется дополнительный эксперимент, и не оставляет эти решения ``на усмотрение после просмотра результатов'' \cite{hooker1995,johnson1999}.

\subsection{Набор публикуемых метрик}

В итоговых таблицах по каждой страте \((family,n,m)\) рекомендуется приводить:
\begin{enumerate}[label=\arabic*.]
  \item выборочное среднее значения \texttt{MINSUM};
  \item стандартное отклонение, минимум и максимум;
  \item \(95\%\)-й доверительный интервал для среднего;
  \item среднее, медиану и при необходимости \(p90\) времени работы;
  \item долю валидных запусков;
  \item относительный разрыв к reference baseline;
  \item для парных сравнений --- среднюю разность, её доверительный интервал, \(p\)-значения paired \(t\)-test и Wilcoxon signed-rank.
\end{enumerate}

Для детерминированных solver'ов по качеству решения достаточно одного запуска на экземпляр, но для измерения времени предпочтительно пакетное измерение: одинаковый solver запускается многократно подряд на одном и том же наборе входов, после чего суммарное время делится на число прогонов. Это уменьшает влияние системного шума при измерении очень быстрых процедур.

\section{Правила интерпретации сравнения с OR-Tools}

Сравнение с OR-Tools в текущем проекте следует трактовать строго как сравнение с \textbf{внешним reference baseline}. Оно полезно для оценки того, насколько собственные solver'ы приближаются к сильному внешнему ориентиру по качеству решения, но не должно описываться как доказательство превосходства по критерию ``качество--время'' без отдельной matched-budget постановки.

В текущем контуре OR-Tools получает фиксированные \(5\) секунд на экземпляр, тогда как собственные solver'ы в среднем работают существенно быстрее. Поэтому формулировка вида ``наш solver лучше OR-Tools, потому что даёт сравнимое качество за меньшее время'' без отдельного эксперимента некорректна. Корректная интерпретация следующая:
\begin{enumerate}[label=\arabic*.]
  \item основной reference-эксперимент отвечает на вопрос о качестве решений относительно внешнего baseline;
  \item отдельный matched-budget эксперимент, если он понадобится в основной главе, должен давать всем solver'ам одинаковый временной бюджет на экземпляр и интерпретироваться отдельно от reference-сводки.
\end{enumerate}

Именно такое разведение уровней сравнения соответствует хорошей практике экспериментальной оценки эвристик и защищает текст курсовой от методически слабых заявлений \cite{hooker1995,johnson1999}.

\section{Итоговый вывод}

Проведённый анализ показывает, что текущая экспериментальная база проекта уже достаточна для двух задач. Во-первых, она подтверждает полезные предварительные содержательные выводы: на uniform-наборе \texttt{lkh-wrapper-v2} является сильнейшим solver'ом по среднему gap к reference baseline и выигрывает у \texttt{grasp} в большинстве страт. Во-вторых, она даёт репозиторию именно тот набор локальных статистических артефактов, который нужен для фиксации защищаемого экспериментального протокола.

Одновременно текущие CSV наглядно показывают, почему старую схему нельзя оставлять в исходном виде как итоговую методику: основной грид пока отражает только uniform-инстансы, а в главном CSV нет независимых повторов стохастических solver'ов по различным \texttt{seed}. Поэтому итоговый текст курсовой должен явно различать уже подтверждённое состояние контура и финальный протокол, к которому проект приведён на уровне конфигурации и pilot-расчётов.

В зафиксированном виде методика принимает следующую окончательную форму: полный грид по четырём семействам инстансов и \(56\) стратам, базовый режим \(N=10, R=20\), подтверждающий режим \(N=20, R=30\), локальный строгий режим \(R=50\), paired-сравнение на одних и тех же экземплярах, paired \(t\)-test и Wilcoxon как взаимодополняющие инструменты, процедура Холма при множественных сравнениях и отдельная трактовка OR-Tools как внешнего reference baseline. Именно в таком виде экспериментальная часть становится статистически замкнутой и пригодной для включения в итоговый текст курсовой работы.

\printbibliography[title={Список литературы}]

\end{refsection}
